\cleardoublepage
\newrefsection
\chapter{文献综述}

\section{背景介绍}

\subsection{多模态大模型与视频大模型发展}

近年来，大规模预训练模型的发展显著推动了视觉理解与跨模态推理技术的进步。早期研究主要集中于视觉-语言统一表示学习，其中 CLIP 模型通过在大规模图文对上进行对比学习，实现视觉特征与语言语义空间的有效对齐，使模型能够在无需任务特定训练的情况下完成多种视觉识别任务，从而为多模态大模型的研究奠定了重要基础\cite{clip2021}。

在此基础上，Flamingo 模型进一步提出在冻结语言模型的基础上引入跨模态注意力机制，使模型能够处理交错的视觉与文本序列输入，并在多种视觉语言任务中实现少样本推理能力，标志着多模态生成式模型架构的初步形成\cite{flamingo2022}。随后，BLIP-2 模型提出利用轻量级的 Q-Former 模块建立视觉编码器与大型语言模型之间的桥接表示，通过冻结视觉编码器和语言模型显著降低训练成本，同时保持较强的跨模态推理能力\cite{blip22023}。

进一步地，LLaVA 与 MiniGPT-4 等研究通过视觉指令微调（visual instruction tuning）策略，使模型能够在自然语言指令驱动下完成复杂视觉理解与推理任务。这类方法通常采用视觉特征投影模块将视觉 token 映射至语言模型特征空间，从而形成当前多模态大模型的主流结构范式\cite{llava2023,minigpt42023}。

随着多模态模型能力不断提升，研究逐渐从静态图像理解扩展至视频理解任务。Video-LLaMA 模型首次实现音频、视频与文本三模态联合建模，通过时序视觉特征编码与跨模态融合机制，使大型语言模型具备视频问答与动态场景理解能力\cite{videollama2023}。Video-LLaVA 在此基础上探索统一视觉表示方式，提高了视频推理性能与模型泛化能力\cite{videollava2023}。最新提出的 Video-LLaMA-2 模型通过引入时空卷积连接器（Spatio-Temporal Convolution connector），进一步强化了视频长时序关系建模能力，标志着视频多模态模型逐渐从简单特征拼接向专用时空结构设计演进\cite{videollama22024}。

此外，LLaVA-Octopus 等工作开始关注多模态融合效率问题，通过设计自适应视觉投影结构，在空间特征与时间特征融合路径之间动态选择，从而提升视觉 token 利用效率并降低推理开销\cite{llavaoctopus2025}。总体而言，多模态大模型的发展经历了从视觉-语言表示学习到统一多模态推理框架，再到面向视频时空理解的体系结构演进过程。

\subsection{视频大模型高效推理问题的提出}

随着视频多模态大模型在智能安防、自动驾驶与虚拟现实等场景中的广泛应用，其推理效率问题逐渐成为制约实际部署的重要因素。一方面，视频理解任务通常涉及长时序帧序列建模，视觉 token 数量随视频长度显著增长，从而导致模型推理计算复杂度迅速上升。相关研究指出，在长视频场景中，上下文窗口限制与高维视觉特征输入会显著影响模型推理效率与响应速度\cite{llavamr2024}。

另一方面，多模态模型通常采用视觉、音频与文本联合推理机制。例如 Video-LLaMA 通过多模态编码器实现跨模态信息融合，但该结构也显著增加了注意力计算规模，使得模型推理延迟明显增加\cite{videollama2023}。此外，Video-LLaMA-2 表明，为增强视频时空关系建模能力，模型需要引入高维特征表示与专用时空模块，从而进一步加剧显存占用问题\cite{videollama22024}。

针对上述挑战，近年来研究开始探索视觉 token 压缩与冗余计算消除策略。例如 LLaVA-MR 提出通过密集帧编码与动态视觉 token 压缩降低长视频推理复杂度，从而提高模型在长时序场景中的推理效率\cite{llavamr2024}。SAISA 进一步分析视觉 token 注意力冗余现象，证明通过减少视觉 token 间自注意力计算，可在几乎不影响模型性能的情况下显著降低推理 FLOPs\cite{saisa2025}。因此，如何在保证视频理解性能的同时降低推理计算开销与显存占用，已成为当前视频多模态模型研究的重要问题。

\section{国内外研究现状}

\subsection{研究方向及进展}

围绕视频多模态模型高效推理问题，国内外学者主要从视觉 token 压缩、时空建模优化、多模态融合结构改进以及长视频推理机制等方向开展研究。

在视觉 token 压缩方面，Token Merging 方法通过合并相似视觉 token 降低 Transformer 模型的输入序列长度，从而减少计算复杂度\cite{tome2023}；DynamicViT 则通过动态剪枝策略在推理过程中选择性保留重要视觉 token，提高模型计算效率\cite{dynamicvit2021}。此外，LLaVA-MR 与 SAISA 等研究将视觉 token 压缩策略扩展至多模态模型推理阶段，从模型结构与注意力机制两个层面降低视觉计算冗余\cite{llavamr2024,saisa2025}。

在时空建模优化方面，TimeSformer 提出将视频 Transformer 的空间注意力与时间注意力分离建模，从而降低视频建模计算成本\cite{timesformer2021}；ViViT 进一步探索纯 Transformer 架构在视频理解任务中的应用，验证了基于视觉 token 的长序列建模能力\cite{vivit2021}。InternVideo 则通过统一视频预训练框架提高视频特征表达能力，为视频多模态模型的发展提供了重要技术支撑\cite{internvideo2023}。

在多模态融合结构方面，Flamingo 与 BLIP-2 分别从跨模态注意力设计与桥接模块构建角度提升视觉与语言特征对齐效率\cite{flamingo2022,blip22023}；LLaVA-Octopus 则通过自适应视觉投影结构减少不必要的跨模态计算，从而提高模型推理效率\cite{llavaoctopus2025}。

在长视频推理机制方面，StreamingLLM 提出滑动窗口记忆机制，实现流式长序列推理能力\cite{streamingllm2023}；LongNet 则通过稀疏扩展注意力结构支持超长上下文建模，为长视频理解提供新的技术路径\cite{longnet2023}。总体来看，现有研究已从模型结构设计与输入压缩两个层面提出多种提升视频多模态模型推理效率的技术路线。

\subsection{存在问题}

尽管上述研究在提升视频多模态模型推理效率方面取得了一定进展，但仍存在一些不足。首先，现有视觉 token 压缩方法多针对静态图像或短视频场景设计，对长时序视频中的时空冗余建模能力仍有限，难以充分利用视频内容的时空稀疏性特征。在实际复杂场景中，大量连续帧往往存在显著的空间与语义相似性，但现有模型仍需对每帧视觉特征进行完整编码与推理计算，导致模型在推理阶段产生明显的重复计算开销，从而影响整体推理效率与系统实时性。其次，多模态融合结构优化研究通常侧重跨模态对齐效率，而对视觉特征冗余与显存占用问题关注不足，在高分辨率视频或长时序输入条件下，视觉 token 数量急剧增长，使模型在显存访问与计算资源分配方面面临更大压力，这在一定程度上限制了视频多模态模型在实际应用场景中的部署能力。再次，长视频推理机制研究虽在上下文建模方面取得一定突破，但相关方法多集中于注意力结构改进或记忆机制设计，尚未与视觉 token 压缩策略形成统一的高效推理框架，缺乏从输入压缩与结构优化协同设计的系统性解决方案。

因此，如何在视频多模态模型结构中引入高效视觉 token 压缩机制，实现面向空间维度与时间维度的联合特征建模，并通过动态计算调度策略减少冗余推理过程，仍是当前研究的重要方向。同时，在保证视频理解性能稳定的前提下探索更具通用性的压缩与融合策略，对于推动视频多模态模型在真实场景中的落地应用具有重要理论意义与工程价值。

\section{研究展望}

综上所述，视频多模态大模型在推理效率方面仍面临计算冗余与显存瓶颈等关键问题。虽然已有研究从视觉 token 压缩、时空建模优化及多模态融合结构改进等方面进行了探索，但针对视频模型时空特征特点的高效视觉 token 压缩方法仍较为缺乏。

因此，将高效视觉特征压缩技术引入视频多模态模型结构中，构建面向长视频场景的时空联合压缩机制具有重要研究意义。本研究拟借鉴 TokenPacker 中提出的高效视觉投影与 token 压缩思想，将其迁移应用于 VideoLLaMA2-7B-16F 模型，在视觉编码器输出与 STC connector 之间引入 TokenPacker 模块，实现视觉 token 的自适应压缩与特征重组织。该方法有望在保持视频理解性能的同时显著降低模型推理计算量与显存占用，从而提升视频多模态模型在复杂场景中的实时推理能力。若该方案能够在较小性能损失下显著减少视觉 token 数量、降低推理耗时和显存占用，则说明 TokenPacker 从图像多模态向视频大模型迁移具有可行性，也为后续进一步加入时间维自适应压缩、关键帧筛选或任务引导压缩奠定基础。

\newpage
\begingroup
    \linespreadsingle{}
    \printbibliography[title={参考文献}]
\endgroup
